% ============================================================
%  curso_python_secundaria.tex
%  Estilo: PREAMBULO_INFOGRAFIA (nuevo look dark-mode / tecnológico)
%  El preámbulo base se importó desde execution/estilo_infografia.py
% ============================================================

\documentclass[11pt,a4paper]{article}

% ── Codificación, fuentes y lenguaje ─────────────────────────────────────────
\usepackage[utf8]{inputenc}
\usepackage[T1]{fontenc}
\usepackage[default,scale=0.95]{sourcesanspro}
\renewcommand{\familydefault}{\sfdefault}
\usepackage[spanish,es-noshorthands,es-tabla]{babel}

% ── Geometría y espaciado ────────────────────────────────────────────────────
\usepackage[top=2.2cm, bottom=2.5cm, left=2.2cm, right=2.2cm, headheight=15pt]{geometry}
\usepackage{setspace}
\setstretch{1.18}
\usepackage{parskip}

% ── Matemáticas y unidades ───────────────────────────────────────────────────
\usepackage{amsmath}
\usepackage{amssymb}
\usepackage{siunitx}

% ── Tablas ───────────────────────────────────────────────────────────────────
\usepackage{booktabs}
\usepackage{tabularx}
\usepackage{array}
\usepackage{longtable}
\usepackage{colortbl}
\usepackage{enumitem}

% ── Gráficos ────────────────────────────────────────────────────────────────
\usepackage{graphicx}

% ── Código Python ────────────────────────────────────────────────────────────
\usepackage{listings}

% ── Colores, cajas e iconos ──────────────────────────────────────────────────
\usepackage[dvipsnames]{xcolor}
\usepackage{tcolorbox}
\tcbuselibrary{skins, breakable}
\usepackage{fontawesome5}

% ── Secciones con barra de color ─────────────────────────────────────────────
\usepackage{titlesec}
\titleformat{\section}
  {\normalfont\LARGE\bfseries\color{azulNoche}}
  {\colorbox{cyanNeon}{\color{fondoOscuro}\thesection}}{0.7em}{}
  [\vspace{2pt}\color{cyanNeon}\rule{\linewidth}{1.8pt}]
\titlespacing*{\section}{0pt}{24pt}{10pt}
\titleformat{\subsection}
  {\normalfont\large\bfseries\color{azulNoche}}
  {\textcolor{cyanNeon}{\thesubsection}}{0.7em}{}
  [\vspace{1pt}\color{grisLinea}\rule{\linewidth}{0.6pt}]
\titlespacing*{\subsection}{0pt}{14pt}{6pt}
\titleformat{\subsubsection}
  {\normalfont\normalsize\bfseries\color{azulNoche}}
  {\textcolor{cyanNeon}{\thesubsubsection}}{0.7em}{}
\titlespacing*{\subsubsection}{0pt}{10pt}{4pt}

% ── Cabeceras y pies de página ───────────────────────────────────────────────
\usepackage{fancyhdr}
\usepackage{lastpage}
\pagestyle{fancy}
\fancyhf{}
\renewcommand{\headrulewidth}{0pt}
\renewcommand{\footrulewidth}{0pt}
\fancyhead[L]{\small\color{azulNoche}\textbf{Curso de Python} --- Prof. César Rodríguez}
\fancyhead[R]{\small\color{grisTexto}\thepage}
\fancyfoot[C]{%
  \begin{minipage}{\linewidth}
    \vspace{2pt}
    \color{cyanNeon}\rule{\linewidth}{1.2pt}\\[2pt]
    \centering\color{grisTexto}\footnotesize
    Página \thepage\ de \pageref{LastPage}
  \end{minipage}%
}

% ── Hiperenlaces ─────────────────────────────────────────────────────────────
\usepackage[colorlinks=true, linkcolor=azulNoche, urlcolor=cyanNeon,
            pdftitle={Curso de Programación en Python para Estudiantes de Secundaria}]{hyperref}
\sloppy
\emergencystretch 3em

% ── Paleta de colores — tecnológico dark-mode ─────────────────────────────────
\definecolor{fondoOscuro}{HTML}{0D1B2A}
\definecolor{verdeTurquesa}{HTML}{004D40}
\definecolor{azulNoche}{HTML}{1B3A6B}
\definecolor{azulMedio}{HTML}{2A5298}
\definecolor{cyanNeon}{HTML}{00C8E0}
\definecolor{verdeSignal}{HTML}{00B887}
\definecolor{naranjaVivo}{HTML}{FF7043}
\definecolor{rojoAlerta}{HTML}{E53935}
\definecolor{grisPapel}{HTML}{F4F6FA}
\definecolor{grisLinea}{HTML}{C8D0E0}
\definecolor{grisTexto}{HTML}{6B7A99}
\definecolor{amarilloNota}{HTML}{FFB300}

% Colores específicos para bloques de código Python
\definecolor{codigoFondo}{HTML}{1E2D3D}
\definecolor{codigoComentario}{HTML}{7A8FA6}
\definecolor{codigoCadena}{HTML}{FF9D5C}
\definecolor{codigoKeyword}{HTML}{00C8E0}
\definecolor{codigoNumero}{HTML}{00E5A0}

% ── Banda de título ───────────────────────────────────────────────────────────
\newcommand{\bandaTitulo}[2]{%
  \begin{tcolorbox}[
    enhanced,
    colback=verdeTurquesa, colframe=verdeTurquesa,
    boxrule=0pt, arc=8pt, outer arc=8pt,
    width=\linewidth,
    top=18pt, bottom=6pt, left=14pt, right=14pt,
    borderline south={3.5pt}{0pt}{cyanNeon},
  ]
    \begin{minipage}[c]{0.07\linewidth}
      \centering
      \textcolor{cyanNeon}{\fontsize{30}{30}\selectfont\faIcon{python}}
    \end{minipage}%
    \hspace{8pt}%
    \begin{minipage}[c]{0.88\linewidth}
      {\fontsize{20}{24}\selectfont\bfseries\color{white}#1}\par\vspace{5pt}%
      \textcolor{cyanNeon!80!white}{\small\faIcon{angle-right}~#2}%
    \end{minipage}
    \vspace{4pt}
  \end{tcolorbox}%
  \vspace{8pt}%
}
\newcommand{\iconotexto}[2]{\textcolor{cyanNeon}{\faIcon{#1}}~#2}

% ── Configuración de listings para Python ────────────────────────────────────
\lstdefinestyle{pythonStyle}{
    language=Python,
    backgroundcolor=\color{codigoFondo},
    basicstyle=\ttfamily\small\color{white},
    keywordstyle=\color{codigoKeyword}\bfseries,
    commentstyle=\color{codigoComentario}\itshape,
    stringstyle=\color{codigoCadena},
    numberstyle=\tiny\color{codigoComentario},
    numbers=left,
    numbersep=10pt,
    showstringspaces=false,
    breaklines=true,
    breakatwhitespace=false,
    tabsize=4,
    frame=none,
    rulecolor=\color{azulNoche},
    captionpos=b,
    aboveskip=6pt,
    belowskip=4pt,
    xleftmargin=14pt,
    framexleftmargin=14pt,
    literate=
      {á}{{\'{a}}}1 {é}{{\'{e}}}1 {í}{{\'{i}}}1 {ó}{{\'{o}}}1 {ú}{{\'{u}}}1
      {Á}{{\'{A}}}1 {É}{{\'{E}}}1 {Í}{{\'{I}}}1 {Ó}{{\'{O}}}1 {Ú}{{\'{U}}}1
      {ñ}{{\~{n}}}1 {Ñ}{{\~{N}}}1 {ü}{{\"{u}}}1 {Ü}{{\"{U}}}1
      {¿}{{?`}}1 {¡}{{!`}}1
}

% Envuelve el lstlisting en una caja oscura redondeada
\lstset{style=pythonStyle}
\BeforeBeginEnvironment{lstlisting}{%
  \begin{tcolorbox}[
    enhanced, breakable,
    arc=6pt, outer arc=6pt,
    colback=codigoFondo, colframe=azulNoche,
    boxrule=1pt,
    drop fuzzy shadow=azulNoche!25!white,
    top=2pt, bottom=2pt, left=0pt, right=0pt,
  ]%
}
\AfterEndEnvironment{lstlisting}{\end{tcolorbox}}

% ── Cajas de contenido del curso ─────────────────────────────────────────────
\newtcolorbox{cajaRecuerda}{
    enhanced, breakable,
    arc=6pt, outer arc=6pt,
    colback=azulNoche!10!grisPapel,
    colframe=azulNoche, boxrule=1pt,
    borderline west={4pt}{0pt}{cyanNeon},
    fonttitle=\bfseries\small\color{white},
    title={\faIcon{info-circle}~Recuerda},
    attach boxed title to top left={yshift=-3mm, xshift=6mm},
    boxed title style={colback=azulNoche, arc=4pt, boxrule=0pt},
    drop fuzzy shadow=azulNoche!25!white,
    left=10pt, right=10pt, top=8pt, bottom=8pt,
}

\newtcolorbox{cajaConcepto}{
    enhanced, breakable,
    arc=6pt, outer arc=6pt,
    colback=verdeSignal!8!white,
    colframe=verdeSignal, boxrule=1pt,
    borderline west={4pt}{0pt}{verdeSignal},
    fonttitle=\bfseries\small\color{white},
    title={\faIcon{lightbulb}~Concepto clave},
    attach boxed title to top left={yshift=-3mm, xshift=6mm},
    boxed title style={colback=verdeSignal!80!black, arc=4pt, boxrule=0pt},
    drop fuzzy shadow=verdeSignal!20!white,
    left=10pt, right=10pt, top=8pt, bottom=8pt,
}

\newtcolorbox{cajaEjemplo}{
    enhanced, breakable,
    arc=6pt, outer arc=6pt,
    colback=cyanNeon!5!white,
    colframe=cyanNeon!70!azulNoche, boxrule=1pt,
    borderline west={4pt}{0pt}{cyanNeon},
    fonttitle=\bfseries\small\color{fondoOscuro},
    title={\faIcon{code}~Ejemplo},
    attach boxed title to top left={yshift=-3mm, xshift=6mm},
    boxed title style={colback=cyanNeon, arc=4pt, boxrule=0pt},
    drop fuzzy shadow=azulNoche!20!white,
    left=10pt, right=10pt, top=8pt, bottom=8pt,
}

\newtcolorbox{cajaEjercicio}{
    enhanced, breakable,
    arc=6pt, outer arc=6pt,
    colback=azulMedio!7!white,
    colframe=azulMedio, boxrule=1pt,
    borderline west={4pt}{0pt}{azulMedio},
    fonttitle=\bfseries\small\color{white},
    title={\faIcon{pencil-alt}~Ejercicio},
    attach boxed title to top left={yshift=-3mm, xshift=6mm},
    boxed title style={colback=azulMedio, arc=4pt, boxrule=0pt},
    drop fuzzy shadow=azulNoche!20!white,
    left=10pt, right=10pt, top=8pt, bottom=8pt,
}

\newtcolorbox{cajaReto}{
    enhanced, breakable,
    arc=6pt, outer arc=6pt,
    colback=naranjaVivo!6!white,
    colframe=naranjaVivo, boxrule=1pt,
    borderline west={4pt}{0pt}{naranjaVivo},
    fonttitle=\bfseries\small\color{white},
    title={\faIcon{fire}~Reto},
    attach boxed title to top left={yshift=-3mm, xshift=6mm},
    boxed title style={colback=naranjaVivo, arc=4pt, boxrule=0pt},
    drop fuzzy shadow=naranjaVivo!20!white,
    left=10pt, right=10pt, top=8pt, bottom=8pt,
}

% ── Indicadores de dificultad ─────────────────────────────────────────────────
\newcommand{\facil}{\textcolor{verdeSignal}{\faIcon{circle}\,\textbf{Fácil}}}
\newcommand{\intermedio}{\textcolor{naranjaVivo}{\faIcon{circle}\,\textbf{Intermedio}}}
\newcommand{\dificil}{\textcolor{rojoAlerta}{\faIcon{circle}\,\textbf{Difícil}}}

% ─────────────────────────────────────────────────────────────────────────────

\begin{document}

% ══ PORTADA ═══════════════════════════════════════════════════════════════════
\bandaTitulo{Curso de Programación en Python}{Para Estudiantes de Secundaria --- Prof. César Rodríguez}

\begin{center}
  \vspace{2pt}
  {\large\color{grisTexto} 12 módulos · teoría, ejemplos y ejercicios · sin requisitos previos}
  \vspace{4pt}
\end{center}

\begin{tcolorbox}[
  enhanced, breakable,
  arc=6pt, outer arc=6pt,
  colback=grisPapel, colframe=grisLinea, boxrule=0.8pt,
  left=12pt, right=12pt, top=8pt, bottom=8pt,
]
\noindent Este curso introduce a los estudiantes de secundaria en el mundo de la programación utilizando el lenguaje Python. A lo largo de 12 módulos se desarrollan los fundamentos de la programación: variables, tipos de datos, estructuras de control, colecciones, funciones y un proyecto final integrador. Cada módulo combina teoría breve, ejemplos resueltos y ejercicios propuestos con distintos niveles de dificultad, de modo que el estudiante aprende haciendo. No se requiere experiencia previa en programación ni conocimientos avanzados de matemáticas.
\end{tcolorbox}

\vspace{6pt}
\tableofcontents
\newpage

% ============================================================
% MÓDULO 1
% ============================================================
\section{Introducción a la Programación}

\subsection{¿Qué es programar?}

Programar consiste en escribir \textbf{instrucciones} para que una computadora realice una tarea. Un \textbf{algoritmo} es la secuencia ordenada de pasos que se necesitan para resolver un problema.

\begin{cajaEjemplo}
Imagina que quieres darle instrucciones a un amigo para preparar un sándwich. Los pasos podrían ser:

\begin{enumerate}
    \item Conseguir pan, queso y jamón.
    \item Cortar el pan por la mitad.
    \item Colocar el queso y el jamón entre las dos mitades.
    \item Servir el sándwich.
\end{enumerate}

Ese conjunto de pasos ordenados es un \textbf{algoritmo}. Una computadora necesita que le demos instrucciones mucho más detalladas, en un lenguaje que ella pueda entender: un \textbf{lenguaje de programación}.
\end{cajaEjemplo}

\subsection{¿Qué es Python?}

Python es un lenguaje de programación creado por Guido van Rossum a finales de los años 80. Algunas de sus características más importantes son:

\begin{itemize}
    \item \textbf{Fácil de leer:} su sintaxis es clara y cercana al lenguaje humano.
    \item \textbf{Fácil de escribir:} se necesitan pocas líneas para hacer mucho.
    \item \textbf{Muy usado:} se emplea en ciencia de datos, inteligencia artificial, robótica, desarrollo web y videojuegos, entre otros campos.
    \item \textbf{Gratis y multiplataforma:} funciona en Windows, Linux y macOS.
\end{itemize}

\begin{cajaRecuerda}
La versión estable de Python es la \textbf{3.14} (año 2026). Todas las computadoras de esta era usan Python 3; Python 2 ya no se usa.
\end{cajaRecuerda}

\subsection{Instalación y primer programa}

\subsubsection{Instalación en Windows}
\begin{enumerate}
    \item Ingresa a \texttt{https://www.python.org/downloads/}.
    \item Descarga el instalador de la última versión (Python 3.14).
    \item Ejecuta el instalador y \textbf{marca la casilla} ``Add Python to PATH''.
    \item Haz clic en ``Install Now''.
\end{enumerate}

\subsubsection{Instalación en Linux}
Python suele venir preinstalado. Para verificarlo, abre una terminal y escribe:

\begin{lstlisting}[style=pythonStyle]
python3 --version
\end{lstlisting}

Si no está instalado, usa el gestor de paquetes de tu distribución:

\begin{lstlisting}[style=pythonStyle]
sudo apt install python3
\end{lstlisting}

\subsubsection{Alternativa en línea}

Si no puedes instalar Python en tu computadora, puedes usar \textbf{Replit} (\texttt{replit.com}), \textbf{Google Colab} (\texttt{colab.research.google.com}) o \textbf{Python Tutor} (\texttt{pythontutor.com}) para practicar sin instalar nada.

\subsection{El primer programa: ``Hola mundo''}

El programa más sencillo que se puede escribir en Python es:

\begin{lstlisting}[style=pythonStyle, caption={Nuestro primer programa}]
print("Hola, mundo!")
\end{lstlisting}

Cuando se ejecuta, el programa muestra en pantalla:

\begin{lstlisting}[style=pythonStyle]
Hola, mundo!
\end{lstlisting}

\subsubsection{¿Cómo se ejecuta?}

\begin{enumerate}
    \item Abre el \textbf{IDLE} (viene instalado con Python) o un editor como Thonny o VS Code.
    \item Escribe el programa anterior en un archivo nuevo.
    \item Guarda el archivo con extensión \texttt{.py} (por ejemplo, \texttt{hola.py}).
    \item Ejecuta el programa con la tecla \texttt{F5} (en IDLE) o con el botón ``Run''.
\end{enumerate}

\begin{cajaConcepto}
\texttt{print()} es una \textbf{función} que muestra texto en la pantalla. El texto va entre comillas simples \texttt{'} o dobles \texttt{"}. Casi todos los programas que escribamos comenzarán usando \texttt{print()}.
\end{cajaConcepto}

\subsubsection{Comentarios en el código}

Un \textbf{comentario} es una nota que escribimos para los humanos que lean el código. La computadora lo ignora. En Python los comentarios empiezan con el símbolo \texttt{\#}:

\begin{lstlisting}[style=pythonStyle, caption={Programa con comentarios}]
# Este es mi primer programa
# Muestra un saludo en pantalla
print("Hola, mundo!")
print("Bienvenidos al curso de Python")
\end{lstlisting}

\subsection{Ejercicios propuestos}

\begin{cajaEjercicio}
\textbf{Ejercicio 1} \facil

Escribe un programa que muestre tu nombre completo en pantalla, por ejemplo:
\begin{lstlisting}[style=pythonStyle]
Mi nombre es Ana López
\end{lstlisting}
\end{cajaEjercicio}

\begin{cajaEjercicio}
\textbf{Ejercicio 2} \facil

Escribe un programa que muestre tres líneas distintas:
\begin{enumerate}
    \item Tu nombre.
    \item Tu ciudad.
    \item Tu hobby favorito.
\end{enumerate}
Cada dato debe aparecer en una línea separada usando tres llamadas a \texttt{print()}.
\end{cajaEjercicio}

\begin{cajaEjercicio}
\textbf{Ejercicio 3} \facil

Escribe un programa que muestre un rectángulo dibujado con asteriscos, como este:
\begin{lstlisting}[style=pythonStyle]
*****
*****
*****
\end{lstlisting}
Ayuda: cada línea del rectángulo es una llamada a \texttt{print()}.
\end{cajaEjercicio}

\begin{cajaEjercicio}
\textbf{Ejercicio 4} \intermedio

Escribe un programa con al menos tres comentarios que explique, paso a paso, qué hace el programa ``Hola mundo''. No es necesario que el programa haga algo nuevo; la idea es practicar la escritura de comentarios claros.
\end{cajaEjercicio}

\begin{cajaReto}
\textbf{Reto 1} \dificil

Investiga qué hace la función \texttt{print()} cuando recibe dos argumentos separados por coma, por ejemplo:
\begin{lstlisting}[style=pythonStyle]
print("Hola", "mundo")
\end{lstlisting}
Escribe un programa que use esta idea para imprimir tu nombre y tu edad en la misma línea, separados por un espacio. Explica con un comentario qué observaste.
\end{cajaReto}

% ============================================================
% MÓDULO 2
% ============================================================
\section{Variables y Tipos de Datos}

\subsection{¿Qué es una variable?}

Una \textbf{variable} es una \textbf{caja} en la memoria de la computadora donde guardamos un valor. A cada caja le ponemos una etiqueta (un \textbf{nombre}) para poder encontrarla después.

\begin{cajaEjemplo}
\begin{lstlisting}[style=pythonStyle, caption={Creando variables}]
edad = 15
nombre = "Ana"
estatura = 1.62

print(edad)
print(nombre)
print(estatura)
\end{lstlisting}

Aquí guardamos:
\begin{itemize}
    \item el número entero 15 en la variable \texttt{edad};
    \item el texto "Ana" en la variable \texttt{nombre};
    \item el número decimal 1.62 en la variable \texttt{estatura}.
\end{itemize}
\end{cajaEjemplo}

\subsection{Reglas para nombrar variables}

\begin{itemize}
    \item Los nombres no pueden empezar con un número: \texttt{2edad} es inválido.
    \item No pueden contener espacios: usa guión bajo \texttt{\_}: \texttt{mi\_edad}.
    \item No pueden usar palabras reservadas de Python (como \texttt{print}, \texttt{if}, \texttt{for}).
    \item Distinguen mayúsculas y minúsculas: \texttt{Edad} y \texttt{edad} son variables diferentes.
    \item Es buena práctica usar nombres descriptivos: \texttt{promedio\_notas} es mejor que \texttt{p}.
\end{itemize}

\begin{cajaRecuerda}
El símbolo \texttt{=} en programación \textbf{no} significa ``igual que''. Significa ``\textbf{asignar}'' o ``guardar en''. La variable de la izquierda recibe el valor de la derecha.
\end{cajaRecuerda}

\subsection{Tipos de datos básicos}

En Python, cada valor tiene un \textbf{tipo}. Los más usados son:

\begin{table}[h]
\centering
\begin{tabularx}{\textwidth}{l l l X}
\toprule
\textbf{Tipo} & \textbf{Nombre} & \textbf{Ejemplo} & \textbf{Descripción} \\
\midrule
Entero & \texttt{int} & \texttt{42} & Números sin parte decimal \\
Decimal & \texttt{float} & \texttt{3.14} & Números con parte decimal \\
Texto & \texttt{str} & \texttt{"Hola"} & Cadenas de caracteres \\
Lógico & \texttt{bool} & \texttt{True} & Solo dos valores: Verdadero o Falso \\
\bottomrule
\end{tabularx}
\end{table}

Para saber el tipo de una variable usamos la función \texttt{type()}:

\begin{lstlisting}[style=pythonStyle, caption={La función type()}]
edad = 15
print(type(edad))          # <class 'int'>

pi = 3.14
print(type(pi))            # <class 'float'>

nombre = "Ana"
print(type(nombre))        # <class 'str'>

es_mayor = True
print(type(es_mayor))      # <class 'bool'>
\end{lstlisting}

\subsection{Cambio de tipo de datos (casting)}

Podemos convertir un valor de un tipo a otro:

\begin{lstlisting}[style=pythonStyle, caption={Conversión de tipos}]
numero_texto = "10"
print(type(numero_texto))          # <class 'str'>

numero = int(numero_texto)         # "10" se convierte en 10
print(type(numero))                # <class 'int'>
print(numero + 5)                  # 15

# De decimal a entero (se truncan los decimales)
precio = 9.99
precio_entero = int(precio)
print(precio_entero)               # 9

# De número a texto
anio = 2026
texto = str(anio)
print(texto + " es el año actual")  # 2026 es el año actual
\end{lstlisting}

\begin{cajaRecuerda}
Las funciones de conversión más usadas son:
\begin{itemize}
    \item \texttt{int(valor)}: convierte a entero.
    \item \texttt{float(valor)}: convierte a decimal.
    \item \texttt{str(valor)}: convierte a texto.
    \item \texttt{bool(valor)}: convierte a lógico.
\end{itemize}
\end{cajaRecuerda}

\subsection{Ejercicios propuestos}

\begin{cajaEjercicio}
\textbf{Ejercicio 5} \facil

Crea tres variables: una con tu edad, una con tu nombre y una con tu estatura. Luego imprime cada una con \texttt{print()}, y además imprime el tipo de cada variable usando \texttt{type()}.
\end{cajaEjercicio}

\begin{cajaEjercicio}
\textbf{Ejercicio 6} \facil

Escribe un programa que calcule y muestre el área de un rectángulo de base 8 y altura 5. Usa dos variables (\texttt{base} y \texttt{altura}) y otra variable \texttt{area} para guardar el resultado. Recuerda que el área se calcula como base $\times$ altura.
\end{cajaEjercicio}

\begin{cajaEjercicio}
\textbf{Ejercicio 7} \intermedio

Encuentra el error en el siguiente programa y corrígelo:
\begin{lstlisting}[style=pythonStyle]
nombre = "Carlos"
edad = 14
print("Hola, me llamo" + nombre + " y tengo" + edad + " años")
\end{lstlisting}
Ayuda: ¿qué tipo de dato espera la suma con \texttt{+}? Ejecuta el programa, observa el mensaje de error y explica la solución en un comentario.
\end{cajaEjercicio}

\begin{cajaEjercicio}
\textbf{Ejercicio 8} \intermedio

Escribe un programa que intercambie el valor de dos variables. Por ejemplo, si \texttt{a = 5} y \texttt{b = 3}, al final debe quedar \texttt{a = 3} y \texttt{b = 5}. Imprime los valores antes y después. Ayuda: usa una variable auxiliar (una ``tercera caja'').
\end{cajaEjercicio}

\begin{cajaReto}
\textbf{Reto 2} \dificil

Escribe un programa que calcule y muestre el perímetro de una circunferencia. Usa una variable \texttt{radio = 7} y el valor de \texttt{pi = 3.14159}. La fórmula del perímetro es $P = 2 \cdot \pi \cdot radio$. Muestra el resultado con un mensaje descriptivo, por ejemplo: \texttt{El perimetro es 43.98226}.
\end{cajaReto}

% ============================================================
% MÓDULO 3
% ============================================================
\section{Entrada de Datos por Teclado}

\subsection{La función input()}

Hasta ahora los programas siempre mostraban los mismos valores. Para que el programa \textbf{interactúe} con el usuario, usamos la función \texttt{input()}, que lee lo que el usuario escribe por teclado.

\begin{cajaEjemplo}
\begin{lstlisting}[style=pythonStyle, caption={Saludando al usuario}]
nombre = input("¿Cómo te llamas? ")
print("Mucho gusto,", nombre)
\end{lstlisting}

Cuando se ejecuta, el programa espera a que el usuario escriba algo y presione \texttt{Enter}. Luego responde:

\begin{lstlisting}[style=pythonStyle]
¿Cómo te llamas? Ana
Mucho gusto, Ana
\end{lstlisting}
\end{cajaEjemplo}

\begin{cajaRecuerda}
\textbf{Importante:} \texttt{input()} siempre devuelve un \textbf{texto} (\texttt{str}), aunque el usuario escriba un número. Si necesitamos operar con números, debemos convertir el valor con \texttt{int()} o \texttt{float()}.
\end{cajaRecuerda}

\subsection{Programa con entrada numérica}

\begin{cajaEjemplo}
\begin{lstlisting}[style=pythonStyle, caption={Suma de dos números}]
numero1 = int(input("Escribe el primer número: "))
numero2 = int(input("Escribe el segundo número: "))

suma = numero1 + numero2
print("La suma es:", suma)
\end{lstlisting}

Nota que \texttt{input()} se \textbf{envuelve} en \texttt{int(...)} para convertir el texto en número. Si no lo hiciéramos, la operación \texttt{numero1 + numero2} sería una concatenación de textos (por ejemplo, \texttt{"5" + "3"} daría \texttt{"53"}).
\end{cajaEjemplo}

\subsection{Ejercicios propuestos}

\begin{cajaEjercicio}
\textbf{Ejercicio 9} \facil

Escribe un programa que pregunte al usuario su nombre y su edad, y luego muestre:
\begin{lstlisting}[style=pythonStyle]
Hola, [nombre]. El próximo año tendrás [edad + 1] años.
\end{lstlisting}
Ayuda: la edad se lee como texto; conviértela con \texttt{int()} para poder sumarle 1.
\end{cajaEjercicio}

\begin{cajaEjercicio}
\textbf{Ejercicio 10} \facil

Escribe un programa que pida dos números enteros y muestre su suma, su resta, su multiplicación y su división. Por cada operación imprime una línea descriptiva, por ejemplo:
\begin{lstlisting}[style=pythonStyle]
5 + 3 = 8
5 - 3 = 2
5 * 3 = 15
5 / 3 = 1.6666666666666667
\end{lstlisting}
\end{cajaEjercicio}

\begin{cajaEjercicio}
\textbf{Ejercicio 11} \intermedio

Escribe un programa que pregunte la base y la altura de un triángulo y calcule su área usando la fórmula $A = \dfrac{b \cdot h}{2}$. Muestra el resultado con un mensaje claro. Ayuda: usa \texttt{float()} para aceptar números con decimales.
\end{cajaEjercicio}

\begin{cajaEjercicio}
\textbf{Ejercicio 12} \intermedio

Escribe un programa que pregunte el precio de un producto y calcule el precio con el 12\% de IVA incluido. El programa debe mostrar el precio original, el monto del IVA y el precio final.
\end{cajaEjercicio}

\begin{cajaReto}
\textbf{Reto 3} \dificil

Escribe un programa que pregunte cuántos segundos tiene una duración y la convierta en \textbf{horas, minutos y segundos}. Por ejemplo, 3661 segundos son 1 hora, 1 minuto y 1 segundo. Ayuda: usa la división entera \texttt{//} y el módulo \texttt{\%} (el residuo de la división). Prueba tu programa con 10000 segundos.
\end{cajaReto}

% ============================================================
% MÓDULO 4
% ============================================================
\section{Operadores y Expresiones}

\subsection{Operadores aritméticos}

\begin{table}[h]
\centering
\begin{tabularx}{\textwidth}{l l l X}
\toprule
\textbf{Operador} & \textbf{Nombre} & \textbf{Ejemplo} & \textbf{Resultado} \\
\midrule
\texttt{+} & Suma & \texttt{7 + 3} & \texttt{10} \\
\texttt{-} & Resta & \texttt{7 - 3} & \texttt{4} \\
\texttt{*} & Multiplicación & \texttt{7 * 3} & \texttt{21} \\
\texttt{/} & División & \texttt{7 / 3} & \texttt{2.333...} \\
\texttt{//} & División entera & \texttt{7 // 3} & \texttt{2} \\
\texttt{\%} & Módulo (residuo) & \texttt{7 \% 3} & \texttt{1} \\
\texttt{**} & Potencia & \texttt{2 ** 5} & \texttt{32} \\
\bottomrule
\end{tabularx}
\end{table}

\begin{cajaConcepto}
El \textbf{módulo} (\texttt{\%}) devuelve el residuo de una división entera. Por ejemplo, \texttt{10 \% 3} es \texttt{1} porque 10 entre 3 es 3 y sobran 1. Es muy útil para saber si un número es par: si \texttt{numero \% 2} es \texttt{0}, el número es par.
\end{cajaConcepto}

\subsection{Operadores de comparación}

Estos operadores comparan dos valores y devuelven un resultado lógico: \texttt{True} (verdadero) o \texttt{False} (falso).

\begin{table}[h]
\centering
\begin{tabularx}{\textwidth}{l l l}
\toprule
\textbf{Operador} & \textbf{Significado} & \textbf{Ejemplo} \\
\midrule
\texttt{==} & Igual que & \texttt{5 == 5} \textrightarrow{} \texttt{True} \\
\texttt{!=} & Diferente de & \texttt{5 != 3} \textrightarrow{} \texttt{True} \\
\texttt{>} & Mayor que & \texttt{5 > 8} \textrightarrow{} \texttt{False} \\
\texttt{<} & Menor que & \texttt{5 < 8} \textrightarrow{} \texttt{True} \\
\texttt{>=} & Mayor o igual que & \texttt{5 >= 5} \textrightarrow{} \texttt{True} \\
\texttt{<=} & Menor o igual que & \texttt{5 <= 4} \textrightarrow{} \texttt{False} \\
\bottomrule
\end{tabularx}
\end{table}

\begin{cajaRecuerda}
No confundas \texttt{=} (asignar un valor a una variable) con \texttt{==} (comparar si dos valores son iguales).
\end{cajaRecuerda}

\subsection{Operadores lógicos}

Permiten combinar condiciones:

\begin{table}[h]
\centering
\begin{tabularx}{\textwidth}{l l X}
\toprule
\textbf{Operador} & \textbf{Significado} & \textbf{Resultado} \\
\midrule
\texttt{and} & Y & Verdadero solo si \textbf{ambas} condiciones son verdaderas \\
\texttt{or} & O & Verdadero si \textbf{al menos una} condición es verdadera \\
\texttt{not} & NO & Invierte el valor: verdadero se convierte en falso y viceversa \\
\bottomrule
\end{tabularx}
\end{table}

\begin{cajaEjemplo}
\begin{lstlisting}[style=pythonStyle, caption={Usando operadores lógicos}]
edad = 16
tiene_permiso = True

# and: las dos deben ser verdaderas
puede_entrar = (edad >= 18) and tiene_permiso
print(puede_entrar)          # False, porque 16 no es >= 18

# or: basta con que una sea verdadera
puede_ver_pelicula = (edad >= 13) or tiene_permiso
print(puede_ver_pelicula)    # True

# not: invierte el valor
no_es_mayor = not (edad >= 18)
print(no_es_mayor)           # True
\end{lstlisting}
\end{cajaEjemplo}

\subsection{Precedencia de operadores}

Cuando una expresión tiene varios operadores, Python los evalúa en un orden establecido (como en matemáticas):

\begin{enumerate}
    \item Paréntesis \texttt{( )}
    \item Potencias \texttt{**}
    \item Multiplicación, división y módulo \texttt{* / // \%}
    \item Suma y resta \texttt{+ -}
    \item Comparaciones \texttt{== != < > <= >=}
    \item Lógicos \texttt{and}, \texttt{or}, \texttt{not}
\end{enumerate}

\begin{cajaRecuerda}
Cuando tengas dudas, \textbf{usa paréntesis}. El código \texttt{(2 + 3) * 4} es más claro que \texttt{2 + 3 * 4} y evita errores de interpretación.
\end{cajaRecuerda}

\subsection{Ejercicios propuestos}

\begin{cajaEjercicio}
\textbf{Ejercicio 13} \facil

Sin ejecutar el programa, escribe en tu cuaderno qué mostrará cada \texttt{print()}. Luego verifica con Python.
\begin{lstlisting}[style=pythonStyle]
print(15 // 4)
print(15 % 4)
print(2 ** 10)
print(10 - 2 * 3)
print((10 - 2) * 3)
print(9 % 2)
\end{lstlisting}
\end{cajaEjercicio}

\begin{cajaEjercicio}
\textbf{Ejercicio 14} \facil

Escribe un programa que le pida al usuario un número y muestre si el número es par o impar. Ayuda: un número es par si \texttt{numero \% 2} es igual a \texttt{0}.
\end{cajaEjercicio}

\begin{cajaEjercicio}
\textbf{Ejercicio 15} \intermedio

Escribe un programa que pida tres notas de un estudiante (de 0 a 20) y calcule el promedio. Luego muestra el promedio con dos decimales. Ayuda: usa la función \texttt{round(promedio, 2)} para redondear.
\end{cajaEjercicio}

\begin{cajaEjercicio}
\textbf{Ejercicio 16} \intermedio

Escribe un programa que pida un número de 4 cifras y muestre sus dígitos por separado. Por ejemplo, si el usuario escribe \texttt{4521}, el programa debe mostrar:
\begin{lstlisting}[style=pythonStyle]
Dígitos: 4 5 2 1
\end{lstlisting}
Ayuda: divide sucesivamente con \texttt{//} y \texttt{\%}. Por ejemplo, \texttt{4521 // 1000 = 4} y \texttt{4521 \% 1000 = 521}.
\end{cajaEjercicio}

\begin{cajaReto}
\textbf{Reto 4} \dificil

Escribe un programa que pida un número entero de tres cifras e indique si es un \textbf{número capicúa} (se lee igual de izquierda a derecha que de derecha a izquierda). Por ejemplo, \texttt{121} es capicúa, pero \texttt{123} no lo es. Ayuda: descompón el número en centenas, decenas y unidades usando \texttt{//} y \texttt{\%}.
\end{cajaReto}

% ============================================================
% MÓDULO 5
% ============================================================
\section{Estructuras Condicionales}

\subsection{La sentencia if}

Un programa no siempre hace lo mismo. A veces necesita \textbf{tomar decisiones}. La estructura \texttt{if} permite ejecutar un bloque de código \textbf{solo si} una condición es verdadera.

\begin{cajaEjemplo}
\begin{lstlisting}[style=pythonStyle, caption={Primer condicional}]
edad = int(input("¿Cuántos años tienes? "))

if edad >= 18:
    print("Eres mayor de edad")
else:
    print("Aún eres menor de edad")
\end{lstlisting}

Si el usuario escribe 20, el programa muestra ``Eres mayor de edad''. Si escribe 14, muestra ``Aún eres menor de edad''.
\end{cajaEjemplo}

\begin{cajaRecuerda}
\begin{itemize}
    \item Después de \texttt{if} (y de \texttt{else}), la línea termina con \textbf{dos puntos} \texttt{:}.
    \item El bloque que se ejecuta debe estar \textbf{sangrado} (indentado) con espacios o tabulaciones. Python usa la sangría para saber qué instrucciones pertenecen a cada bloque.
    \item Por convención se usan \textbf{4 espacios} de sangría.
\end{itemize}
\end{cajaRecuerda}

\subsection{La sentencia elif}

Cuando hay más de dos opciones, usamos \texttt{elif} (contracción de ``else if'').

\begin{cajaEjemplo}
\begin{lstlisting}[style=pythonStyle, caption={Varias opciones con elif}]
nota = int(input("Escribe tu nota (0 a 20): "))

if nota >= 18:
    print("Sobresaliente")
elif nota >= 15:
    print("Notable")
elif nota >= 11:
    print("Aprobado")
else:
    print("Reprobado")
\end{lstlisting}

Python evalúa las condiciones \textbf{en orden}: si la primera es verdadera, ejecuta su bloque y \textbf{se salta las demás}.
\end{cajaEjemplo}

\subsection{Condiciones anidadas}

Dentro de un bloque \texttt{if} podemos escribir otro \texttt{if}. Esto se llama \textbf{anidamiento}.

\begin{cajaEjemplo}
\begin{lstlisting}[style=pythonStyle, caption={Condicionales anidados}]
edad = int(input("¿Cuántos años tienes? "))

if edad >= 12:
    if edad >= 18:
        print("Eres un adulto")
    else:
        print("Eres un adolescente")
else:
    print("Eres un niño o niña")
\end{lstlisting}
\end{cajaEjemplo}

\subsection{Ejercicios propuestos}

\begin{cajaEjercicio}
\textbf{Ejercicio 17} \facil

Escribe un programa que pida un número y diga si es \textbf{positivo}, \textbf{negativo} o \textbf{cero}.
\end{cajaEjercicio}

\begin{cajaEjercicio}
\textbf{Ejercicio 18} \facil

Escribe un programa que pida un número entero y muestre si es \textbf{par} o \textbf{impar}, y además si es \textbf{mayor que 100} o \textbf{menor o igual que 100}. Usa condiciones anidadas.
\end{cajaEjercicio}

\begin{cajaEjercicio}
\textbf{Ejercicio 19} \intermedio

Escribe un programa que pida dos números y muestre cuál es el \textbf{mayor}. Si son iguales, debe mostrar el mensaje ``Los números son iguales''.
\end{cajaEjercicio}

\begin{cajaEjercicio}
\textbf{Ejercicio 20} \intermedio

Escribe un programa que pida tres números y los muestre \textbf{en orden ascendente} (del menor al mayor). Ayuda: compara los tres números usando \texttt{if} y \texttt{elif}, y ordena las condiciones de menor a mayor.
\end{cajaEjercicio}

\begin{cajaReto}
\textbf{Reto 5} \dificil

Crea un pequeño \textbf{asistente de clasificación de edades}:
\begin{itemize}
    \item De 0 a 2 años: ``Bebé''.
    \item De 3 a 11 años: ``Niño(a)''.
    \item De 12 a 17 años: ``Adolescente''.
    \item De 18 a 59 años: ``Adulto''.
    \item 60 o más: ``Adulto mayor''.
    \item Si es menor que 0, muestra ``Edad inválida''.
\end{itemize}
Usa \texttt{and} para combinar los límites inferiores y superiores de cada rango.
\end{cajaReto}

% ============================================================
% MÓDULO 6
% ============================================================
\section{Bucles: for y while}

\subsection{¿Por qué necesitamos bucles?}

Un \textbf{bucle} (o \textbf{laço de repetición}) permite ejecutar un bloque de código \textbf{muchas veces} sin escribirlo muchas veces. Por ejemplo, para mostrar los números del 1 al 10 no es necesario escribir diez \texttt{print()}.

\subsection{El bucle for}

El bucle \texttt{for} recorre una secuencia de elementos (por ejemplo, un rango de números) y ejecuta el bloque una vez por cada elemento.

\begin{cajaEjemplo}
\begin{lstlisting}[style=pythonStyle, caption={Mostrar números del 1 al 5}]
for numero in range(1, 6):
    print(numero)
\end{lstlisting}

Salida:
\begin{lstlisting}[style=pythonStyle]
1
2
3
4
5
\end{lstlisting}

Nota que \texttt{range(1, 6)} genera los números del 1 al 5: el \textbf{inicio es inclusivo} y el \textbf{final es exclusivo} (no se incluye el 6).
\end{cajaEjemplo}

\begin{cajaRecuerda}
\begin{itemize}
    \item \texttt{range(n)} genera números de \texttt{0} a \texttt{n-1}.
    \item \texttt{range(inicio, fin)} genera de \texttt{inicio} a \texttt{fin-1}.
    \item \texttt{range(inicio, fin, paso)} genera saltando de \texttt{paso} en \texttt{paso}.
\end{itemize}
\end{cajaRecuerda}

\subsection{El bucle while}

El bucle \texttt{while} repite un bloque \textbf{mientras} una condición sea verdadera. Es útil cuando no sabemos de antemano cuántas veces se repetirá.

\begin{cajaEjemplo}
\begin{lstlisting}[style=pythonStyle, caption={Contar hasta 5 con while}]
contador = 1

while contador <= 5:
    print(contador)
    contador = contador + 1
\end{lstlisting}

El programa muestra los números del 1 al 5. La línea \texttt{contador = contador + 1} es \textbf{esencial}: sin ella, la condición nunca cambiaría y el bucle se repetiría \textbf{infinitamente}.
\end{cajaEjemplo}

\begin{cajaRecuerda}
\begin{itemize}
    \item \textbf{Cuidado con los bucles infinitos:} si la condición de \texttt{while} nunca se vuelve falsa, el programa nunca termina.
    \item Una forma abreviada de \texttt{contador = contador + 1} es \texttt{contador += 1}.
    \item La sentencia \texttt{break} \textbf{sale} del bucle inmediatamente.
    \item La sentencia \texttt{continue} \textbf{salta a la siguiente} iteración.
\end{itemize}
\end{cajaRecuerda}

\subsection{Acumuladores}

Un patrón muy común es usar una variable para \textbf{acumular} la suma de varios valores.

\begin{cajaEjemplo}
\begin{lstlisting}[style=pythonStyle, caption={Sumar los primeros 100 números}]
suma = 0

for numero in range(1, 101):
    suma = suma + numero

print("La suma del 1 al 100 es:", suma)   # 5050
\end{lstlisting}

La variable \texttt{suma} inicia en 0 y en cada vuelta recibe su valor anterior más el número actual.
\end{cajaEjemplo}

\subsection{Ejercicios propuestos}

\begin{cajaEjercicio}
\textbf{Ejercicio 21} \facil

Escribe un programa con \texttt{for} que muestre los números \textbf{pares} del 2 al 20. Ayuda: usa \texttt{range(2, 21, 2)}.
\end{cajaEjercicio}

\begin{cajaEjercicio}
\textbf{Ejercicio 22} \facil

Escribe un programa que muestre la \textbf{tabla de multiplicar} de un número que pide al usuario. Por ejemplo, si el usuario escribe 7, debe mostrar:
\begin{lstlisting}[style=pythonStyle]
7 x 1 = 7
7 x 2 = 14
...
7 x 10 = 70
\end{lstlisting}
\end{cajaEjercicio}

\begin{cajaEjercicio}
\textbf{Ejercicio 23} \intermedio

Escribe un programa que pida al usuario \textbf{5 calificaciones} (una por una) y calcule su promedio. Usa un bucle \texttt{for} y un acumulador. Al final muestra el promedio con dos decimales.
\end{cajaEjercicio}

\begin{cajaEjercicio}
\textbf{Ejercicio 24} \intermedio

Escribe un programa que cuente de 10 hacia atrás hasta 1 (con \texttt{for} y \texttt{range}) y que, al terminar, muestre el mensaje ``¡Despegue!''.
\end{cajaEjercicio}

\begin{cajaEjercicio}
\textbf{Ejercicio 25} \intermedio

Escribe un programa con \texttt{while} que pida al usuario números enteros y los sume, pero que \textbf{se detenga} cuando el usuario escriba 0. Al final, muestra la suma de todos los números escritos.
\end{cajaEjercicio}

\begin{cajaReto}
\textbf{Reto 6} \dificil

Escribe un programa que calcule el \textbf{factorial} de un número. El factorial de $n$ (escrito $n!$) es el producto de todos los enteros desde 1 hasta $n$. Por ejemplo:
\begin{itemize}
    \item $5! = 5 \cdot 4 \cdot 3 \cdot 2 \cdot 1 = 120$
    \item $7! = 5040$
\end{itemize}
Pide el número al usuario y muestra el resultado. Ayuda: usa un acumulador que inicia en 1 y un bucle que vaya de 1 hasta el número.
\end{cajaReto}

% ============================================================
% MÓDULO 7
% ============================================================
\section{Cadenas de Texto}

\subsection{¿Qué es una cadena de texto?}

Una \textbf{cadena} (\texttt{str}) es una secuencia de caracteres. Podemos acceder a cada carácter mediante su \textbf{índice}. Los índices empiezan en \textbf{0}.

\begin{cajaEjemplo}
\begin{lstlisting}[style=pythonStyle, caption={Acceder a caracteres}]
palabra = "Python"

print(palabra[0])    # P
print(palabra[1])    # y
print(palabra[5])    # n
print(len(palabra))  # 6 (cantidad de caracteres)
\end{lstlisting}

Los índices negativos cuentan desde el final:
\begin{lstlisting}[style=pythonStyle]
print(palabra[-1])   # n (el último)
print(palabra[-2])   # o (el penúltimo)
\end{lstlisting}
\end{cajaEjemplo}

\subsection{Rebanadas (slicing)}

Con los dos puntos \texttt{:} podemos extraer una \textbf{porción} de la cadena.

\begin{cajaEjemplo}
\begin{lstlisting}[style=pythonStyle, caption={Rebanadas}]
texto = "Hola mundo"

print(texto[0:4])    # Hola   (del 0 al 3)
print(texto[5:10])   # mundo  (del 5 al 9)
print(texto[:4])     # Hola   (del inicio al 3)
print(texto[5:])     # mundo  (del 5 al final)
print(texto[:])      # Hola mundo (toda la cadena)
\end{lstlisting}

\texttt{cadena[inicio:fin]} toma desde \texttt{inicio} hasta \texttt{fin-1}.
\end{cajaEjemplo}

\subsection{Métodos útiles de cadenas}

\begin{table}[h]
\centering
\begin{tabularx}{\textwidth}{l X}
\toprule
\textbf{Método} & \textbf{Qué hace} \\
\midrule
\texttt{cadena.upper()} & Convierte todo a mayúsculas \\
\texttt{cadena.lower()} & Convierte todo a minúsculas \\
\texttt{cadena.capitalize()} & Mayúscula inicial, resto en minúsculas \\
\texttt{cadena.strip()} & Elimina espacios al inicio y al final \\
\texttt{cadena.count("x")} & Cuenta cuántas veces aparece \texttt{x} \\
\texttt{cadena.find("x")} & Devuelve el índice de la primera aparición de \texttt{x} \\
\texttt{cadena.replace("a","b")} & Reemplaza todas las \texttt{a} por \texttt{b} \\
\texttt{cadena.split(" ") } & Divide la cadena en una lista según el separador \\
\texttt{"-".join(lista)} & Une los elementos de una lista con el separador \\
\bottomrule
\end{tabularx}
\end{table}

\begin{cajaEjemplo}
\begin{lstlisting}[style=pythonStyle, caption={Usando métodos}]
mensaje = "  hola mundo  "

print(mensaje.upper())        # "  HOLA MUNDO  "
print(mensaje.strip())        # "hola mundo"
print(mensaje.strip().upper())# "HOLA MUNDO"

texto = "ana,pedro,luis"
nombres = texto.split(",")
print(nombres)                # ['ana', 'pedro', 'luis']
\end{lstlisting}

Los métodos se pueden \textbf{encadenar}: primero \texttt{strip()} y luego \texttt{upper()}.
\end{cajaEjemplo}

\subsection{Ejercicios propuestos}

\begin{cajaEjercicio}
\textbf{Ejercicio 26} \facil

Escribe un programa que pida una palabra y muestre:
\begin{itemize}
    \item La palabra en mayúsculas.
    \item La palabra en minúsculas.
    \item El número de letras que tiene.
    \item Su primera y última letra.
\end{itemize}
\end{cajaEjercicio}

\begin{cajaEjercicio}
\textbf{Ejercicio 27} \facil

Escribe un programa que pida el nombre completo de una persona y lo muestre con \textbf{mayúscula inicial en cada palabra}. Ayuda: investiga el método \texttt{title()}.
\end{cajaEjercicio}

\begin{cajaEjercicio}
\textbf{Ejercicio 28} \intermedio

Escribe un programa que pida una frase y cuente cuántas veces aparece la letra \texttt{a} (sin importar si es mayúscula o minúscula). Ayuda: convierte la frase a minúsculas con \texttt{lower()} antes de contar.
\end{cajaEjercicio}

\begin{cajaEjercicio}
\textbf{Ejercicio 29} \intermedio

Escribe un programa que pida una palabra e indique si es un \textbf{palíndromo} (se lee igual al derecho y al revés). Ejemplos: \texttt{reconocer}, \texttt{radar}, \texttt{ana}. Ayuda: compara la palabra con su inversa usando \texttt{palabra[::-1]}.
\end{cajaEjercicio}

\begin{cajaEjercicio}
\textbf{Ejercicio 30} \intermedio

Escribe un programa que pida una frase y la muestre con las palabras \textbf{en orden inverso}. Por ejemplo, la frase ``me gusta programar'' debe mostrar ``programar gusta me''. Ayuda: usa \texttt{split()} y \texttt{join()}.
\end{cajaEjercicio}

\begin{cajaReto}
\textbf{Reto 7} \dificil

Escribe un programa que \textbf{cifre} un mensaje con el \textbf{cifrado César}: cada letra se desplaza 3 posiciones en el alfabeto. Por ejemplo, la \texttt{a} se convierte en \texttt{d}, la \texttt{b} en \texttt{e}, y así sucesivamente (la \texttt{z} vuelve a \texttt{a}). Pide al usuario una frase y muestra la versión cifrada. Ayuda: investiga las funciones \texttt{ord()} y \texttt{chr()}, que convierten caracteres en números y viceversa.
\end{cajaReto}

% ============================================================
% MÓDULO 8
% ============================================================
\section{Listas}

\subsection{¿Qué es una lista?}

Una \textbf{lista} es una colección \textbf{ordenada} y \textbf{modificable} de valores. Se crea con corchetes \texttt{[ ]} y sus elementos se separan con comas.

\begin{cajaEjemplo}
\begin{lstlisting}[style=pythonStyle, caption={Creando y usando listas}]
frutas = ["manzana", "pera", "uva", "mango"]

print(frutas)           # ['manzana', 'pera', 'uva', 'mango']
print(frutas[0])        # manzana
print(frutas[2])        # uva
print(len(frutas))      # 4

# Las listas son modificables
frutas[1] = "banana"
print(frutas)           # ['manzana', 'banana', 'uva', 'mango']

# Agregar un elemento al final
frutas.append("kiwi")
print(frutas)           # ['manzana', 'banana', 'uva', 'mango', 'kiwi']

# Quitar el último elemento
ultimo = frutas.pop()
print(ultimo)           # kiwi
print(frutas)
\end{lstlisting}
\end{cajaEjemplo}

\subsection{Métodos y funciones útiles para listas}

\begin{table}[h]
\centering
\begin{tabularx}{\textwidth}{l X}
\toprule
\textbf{Operación} & \textbf{Qué hace} \\
\midrule
\texttt{lista.append(x)} & Agrega \texttt{x} al final \\
\texttt{lista.insert(i, x)} & Inserta \texttt{x} en la posición \texttt{i} \\
\texttt{lista.remove(x)} & Elimina la primera aparición de \texttt{x} \\
\texttt{lista.pop()} & Elimina y devuelve el último elemento \\
\texttt{lista.sort()} & Ordena la lista de menor a mayor (en su lugar) \\
\texttt{sorted(lista)} & Devuelve una copia ordenada sin modificar la original \\
\texttt{sum(lista)} & Suma todos los elementos (si son números) \\
\texttt{min(lista)} & Devuelve el elemento menor \\
\texttt{max(lista)} & Devuelve el elemento mayor \\
\texttt{lista.count(x)} & Cuenta cuántas veces aparece \texttt{x} \\
\texttt{x in lista} & Devuelve \texttt{True} si \texttt{x} está en la lista \\
\bottomrule
\end{tabularx}
\end{table}

\subsection{Recorrer una lista}

Podemos recorrer todos los elementos de una lista con \texttt{for}:

\begin{cajaEjemplo}
\begin{lstlisting}[style=pythonStyle, caption={Recorriendo una lista}]
nombres = ["Ana", "Luis", "Marta", "Pedro"]

for nombre in nombres:
    print("Hola,", nombre)

# También podemos acceder por índice con range
for i in range(len(nombres)):
    print(i, "->", nombres[i])
\end{lstlisting}
\end{cajaEjemplo}

\subsection{Ejercicios propuestos}

\begin{cajaEjercicio}
\textbf{Ejercicio 31} \facil

Crea una lista con 5 colores. Luego:
\begin{enumerate}
    \item Muestra el primer y el último color.
    \item Agrega un sexto color con \texttt{append()}.
    \item Reemplaza el segundo color por otro.
    \item Muestra la lista final y su tamaño con \texttt{len()}.
\end{enumerate}
\end{cajaEjercicio}

\begin{cajaEjercicio}
\textbf{Ejercicio 32} \facil

Escribe un programa que pida al usuario \textbf{5 números} (uno por uno), los guarde en una lista y luego muestre la suma y el promedio. Ayuda: usa \texttt{append()} dentro de un bucle \texttt{for}.
\end{cajaEjercicio}

\begin{cajaEjercicio}
\textbf{Ejercicio 33} \intermedio

Escribe un programa que pida al usuario \textbf{3 notas} (de 0 a 20), las guarde en una lista y muestre:
\begin{itemize}
    \item El promedio.
    \item La nota más alta (\texttt{max}).
    \item La nota más baja (\texttt{min}).
    \item Las notas ordenadas de mayor a menor.
\end{itemize}
\end{cajaEjercicio}

\begin{cajaEjercicio}
\textbf{Ejercicio 34} \intermedio

Escribe un programa que, a partir de una lista de números, cuente cuántos son \textbf{pares} y cuántos son \textbf{impares}. Por ejemplo, con la lista \texttt{[4, 7, 2, 9, 10, 5]}, debe mostrar ``3 pares y 3 impares''.
\end{cajaEjercicio}

\begin{cajaEjercicio}
\textbf{Ejercicio 35} \intermedio

Escribe un programa que pida nombres y los guarde en una lista hasta que el usuario escriba \texttt{"fin"}. Luego muestra:
\begin{itemize}
    \item Cuántos nombres se ingresaron.
    \item El nombre más largo.
    \item La lista ordenada alfabéticamente.
\end{itemize}
\end{cajaEjercicio}

\begin{cajaReto}
\textbf{Reto 8} \dificil

Escribe un programa que genere la \textbf{secuencia de Fibonacci} hasta un límite indicado por el usuario. En la secuencia de Fibonacci cada número es la suma de los dos anteriores: 0, 1, 1, 2, 3, 5, 8, 13, 21, \ldots{} Guarda los primeros términos en una lista y muéstrala. Ayuda: comienza con \texttt{[0, 1]} y ve agregando la suma de los dos últimos.
\end{cajaReto}

% ============================================================
% MÓDULO 9
% ============================================================
\section{Tuplas y Diccionarios}

\subsection{Tuplas}

Una \textbf{tupla} es como una lista, pero \textbf{no se puede modificar} (es \textbf{inmutable}). Se crea con paréntesis \texttt{( )}.

\begin{cajaEjemplo}
\begin{lstlisting}[style=pythonStyle, caption={Creando una tupla}]
punto = (3, 5)
print(punto[0])        # 3
print(punto[1])        # 5

# Esto da ERROR: las tuplas no se pueden modificar
# punto[0] = 10

# Una tupla con un solo elemento necesita una coma
unico = (7,)
print(unico)
\end{lstlisting}

\textbf{Desempaquetado:} podemos asignar cada elemento de una tupla a una variable:
\begin{lstlisting}[style=pythonStyle]
coordenadas = (4, 7)
x, y = coordenadas
print("x =", x)   # x = 4
print("y =", y)   # y = 7
\end{lstlisting}
\end{cajaEjemplo}

\subsection{Diccionarios}

Un \textbf{diccionario} guarda pares de \textbf{clave: valor}. Se crea con llaves \texttt{\{ \}} y permite buscar un valor a partir de su clave.

\begin{cajaEjemplo}
\begin{lstlisting}[style=pythonStyle, caption={Creando un diccionario}]
estudiante = {
    "nombre": "Ana López",
    "edad": 15,
    "curso": "3ro Secundaria"
}

print(estudiante["nombre"])     # Ana López
print(estudiante["edad"])       # 15

# Modificar un valor
estudiante["edad"] = 16
print(estudiante["edad"])       # 16

# Agregar una nueva clave
estudiante["ciudad"] = "Sucre"
print(estudiante)
\end{lstlisting}
\end{cajaEjemplo}

\subsection{Métodos útiles de diccionarios}

\begin{table}[h]
\centering
\begin{tabularx}{\textwidth}{l X}
\toprule
\textbf{Operación} & \textbf{Qué hace} \\
\midrule
\texttt{dic["clave"]} & Devuelve el valor de la clave (error si no existe) \\
\texttt{dic.get("clave")} & Devuelve el valor o \texttt{None} si no existe \\
\texttt{dic.get("clave", x)} & Devuelve \texttt{x} si la clave no existe \\
\texttt{dic.keys()} & Devuelve todas las claves \\
\texttt{dic.values()} & Devuelve todos los valores \\
\texttt{dic.items()} & Devuelve pares (clave, valor) \\
\texttt{"clave" in dic} & Devuelve \texttt{True} si la clave existe \\
\texttt{dic["clave"] = v} & Agrega o modifica el valor de la clave \\
\bottomrule
\end{tabularx}
\end{table}

\begin{cajaEjemplo}
\begin{lstlisting}[style=pythonStyle, caption={Recorriendo un diccionario}]
capitales = {
    "Bolivia": "La Paz",
    "Argentina": "Buenos Aires",
    "Colombia": "Bogotá",
    "Perú": "Lima"
}

for pais, capital in capitales.items():
    print("La capital de", pais, "es", capital)
\end{lstlisting}
\end{cajaEjemplo}

\subsection{Ejercicios propuestos}

\begin{cajaEjercicio}
\textbf{Ejercicio 36} \facil

Crea un diccionario con los datos de un libro: título, autor y año de publicación. Luego:
\begin{enumerate}
    \item Muestra el título.
    \item Modifica el año.
    \item Agrega la clave \texttt{género}.
    \item Muestra todas las claves y todos los valores.
\end{enumerate}
\end{cajaEjercicio}

\begin{cajaEjercicio}
\textbf{Ejercicio 37} \facil

Crea una tupla con las 4 estaciones del año. Muestra la primera y la última estación, y luego la cantidad de elementos con \texttt{len()}. Escribe un comentario explicando por qué no podemos modificar el primer elemento.
\end{cajaEjercicio}

\begin{cajaEjercicio}
\textbf{Ejercicio 38} \intermedio

Escribe un programa que pida al usuario 3 países y su capital (un país a la vez), los guarde en un diccionario y luego muestre las capitales de todos los países ingresados.
\end{cajaEjercicio}

\begin{cajaEjercicio}
\textbf{Ejercicio 39} \intermedio

Dado el diccionario de calificaciones \texttt{\{"Ana": 18, "Luis": 14, "Marta": 20, "Pedro": 11\}}, escribe un programa que:
\begin{itemize}
    \item Muestre el promedio de las calificaciones.
    \item Muestre quién tiene la calificación más alta.
    \item Muestre quién tiene la más baja.
\end{itemize}
Ayuda: usa \texttt{dic.values()}, \texttt{max()} y \texttt{min()}. Para saber quién tiene la nota, busca el estudiante cuya calificación coincida.
\end{cajaEjercicio}

\begin{cajaEjercicio}
\textbf{Ejercicio 40} \intermedio

Escribe un programa que \textbf{cuente las letras} de una palabra usando un diccionario. Por ejemplo, para la palabra \texttt{"banana"}, el resultado debe ser:
\begin{lstlisting}[style=pythonStyle]
{'b': 1, 'a': 3, 'n': 2}
\end{lstlisting}
Ayuda: recorre cada letra; si la letra ya está en el diccionario, suma 1 a su valor; si no, la agrega con valor 1.
\end{cajaEjercicio}

\begin{cajaReto}
\textbf{Reto 9} \dificil

Escribe un programa que gestione una \textbf{agenda telefónica} con un diccionario. El programa debe:
\begin{enumerate}
    \item Permitir agregar contactos (nombre y número).
    \item Permitir buscar el número de un contacto.
    \item Permitir eliminar un contacto.
    \item Mostrar todos los contactos.
\end{enumerate}
El menú se muestra en un bucle \texttt{while} y se sale con la opción ``salir''. Usa \texttt{input()} para elegir la opción.
\end{cajaReto}

% ============================================================
% MÓDULO 10
% ============================================================
\section{Funciones}

\subsection{¿Qué es una función?}

Una \textbf{función} es un bloque de código \textbf{reutilizable} que realiza una tarea concreta. Nos permite escribir el código una sola vez y usarlo muchas veces. Ya hemos usado funciones como \texttt{print()}, \texttt{input()}, \texttt{len()} y \texttt{int()}; ahora aprenderemos a crear las nuestras.

\subsection{Definir y llamar funciones}

Para crear una función usamos la palabra clave \texttt{def}:

\begin{cajaEjemplo}
\begin{lstlisting}[style=pythonStyle, caption={Mi primera función}]
def saludar():
    print("¡Hola! Bienvenido al curso")

# Llamar a la función
saludar()
saludar()
saludar()
\end{lstlisting}

El programa muestra el mensaje tres veces. La definición \textbf{no ejecuta} el código; solo se ejecuta cuando \textbf{llamamos} a la función.
\end{cajaEjemplo}

\subsection{Parámetros y argumentos}

Una función puede recibir \textbf{parámetros} (datos de entrada) para trabajar con ellos.

\begin{cajaEjemplo}
\begin{lstlisting}[style=pythonStyle, caption={Función con parámetros}]
def saludar_persona(nombre):
    print("Hola,", nombre)

def sumar(a, b):
    resultado = a + b
    return resultado

saludar_persona("Ana")
saludar_persona("Luis")

# La función sumar devuelve un valor con return
total = sumar(5, 3)
print("La suma es:", total)
\end{lstlisting}

\begin{itemize}
    \item \texttt{nombre}, \texttt{a} y \texttt{b} son \textbf{parámetros}: variables locales de la función.
    \item \texttt{return} devuelve un resultado para que el programa que llamó a la función pueda usarlo.
\end{itemize}
\end{cajaEjemplo}

\begin{cajaRecuerda}
\textbf{Valores por defecto:} un parámetro puede tener un valor inicial que se usa si la función se llama sin ese argumento:
\begin{lstlisting}[style=pythonStyle]
def saludar(nombre, saludo="Hola"):
    print(saludo + ", " + nombre)

saludar("Ana")              # Hola, Ana
saludar("Luis", "Buenos días")  # Buenos días, Luis
\end{lstlisting}
\end{cajaRecuerda}

\subsection{Ámbito de las variables}

Las variables creadas dentro de una función son \textbf{locales}: solo existen dentro de la función. Las variables creadas fuera son \textbf{globales} y son visibles en todo el programa.

\begin{cajaEjemplo}
\begin{lstlisting}[style=pythonStyle, caption={Variables locales y globales}]
mensaje = "Hola desde fuera"   # variable global

def cambiar_mensaje():
    mensaje_local = "Hola desde dentro"  # variable local
    print(mensaje_local)
    print(mensaje)   # la global también es visible

cambiar_mensaje()
print(mensaje)       # funciona, es global
# print(mensaje_local)  # ERROR: no existe fuera de la función
\end{lstlisting}
\end{cajaEjemplo}

\subsection{Ejercicios propuestos}

\begin{cajaEjercicio}
\textbf{Ejercicio 41} \facil

Escribe una función llamada \texttt{dibujar\_linea()} que muestre una línea de 20 asteriscos. Luego llámala 3 veces para dibujar tres líneas.
\end{cajaEjercicio}

\begin{cajaEjercicio}
\textbf{Ejercicio 42} \facil

Escribe una función \texttt{cuadrado(n)} que reciba un número y devuelva su cuadrado. Pide un número al usuario y muestra el resultado.
\end{cajaEjercicio}

\begin{cajaEjercicio}
\textbf{Ejercicio 43} \intermedio

Escribe una función \texttt{es\_par(n)} que devuelva \texttt{True} si el número es par y \texttt{False} si es impar. Luego escribe un programa que pida 5 números y cuente cuántos son pares usando la función.
\end{cajaEjercicio}

\begin{cajaEjercicio}
\textbf{Ejercicio 44} \intermedio

Escribe una función \texttt{promedio(notas)} que reciba una lista de notas y devuelva su promedio. Luego pide al usuario 4 notas, guárdalas en una lista y muestra el promedio usando la función.
\end{cajaEjercicio}

\begin{cajaEjercicio}
\textbf{Ejercicio 45} \intermedio

Escribe una función \texttt{invertir\_texto(texto)} que devuelva el texto invertido (por ejemplo, ``hola'' se convierte en ``aloh''). Luego escribe un programa que pida una frase, la invierta con la función y muestre el resultado.
\end{cajaEjercicio}

\begin{cajaReto}
\textbf{Reto 10} \dificil

Crea una \textbf{calculadora con funciones}. Define las funciones \texttt{sumar(a, b)}, \texttt{restar(a, b)}, \texttt{multiplicar(a, b)} y \texttt{dividir(a, b)}. Luego escribe un programa que muestre un menú con las 4 operaciones, pida dos números y la operación elegida, y muestre el resultado. Usa un bucle \texttt{while} para que el menú se repita hasta que el usuario elija ``salir''.
\end{cajaReto}

% ============================================================
% MÓDULO 11
% ============================================================
\section{Biblioteca Estándar: math y random}

\subsection{¿Qué es la biblioteca estándar?}

Python incluye cientos de \textbf{módulos} con funciones listas para usar. Un módulo es un archivo con funciones y herramientas relacionadas. Para usarlo, lo importamos con la instrucción \texttt{import}.

\subsection{El módulo math}

Proporciona funciones matemáticas avanzadas:

\begin{cajaEjemplo}
\begin{lstlisting}[style=pythonStyle, caption={Usando math}]
import math

print(math.sqrt(144))        # 12.0  (raíz cuadrada)
print(math.pi)               # 3.141592653589793
print(math.ceil(4.2))        # 5  (redondea hacia arriba)
print(math.floor(4.8))       # 4  (redondea hacia abajo)
print(math.pow(2, 10))       # 1024.0  (potencia)
print(math.factorial(5))     # 120
print(math.fabs(-7))         # 7.0  (valor absoluto)
\end{lstlisting}
\end{cajaEjemplo}

\subsection{El módulo random}

Permite generar \textbf{números aleatorios}, muy útil para juegos:

\begin{cajaEjemplo}
\begin{lstlisting}[style=pythonStyle, caption={Usando random}]
import random

print(random.randint(1, 6))    # número entero entre 1 y 6 (dado)
print(random.random())         # número decimal entre 0 y 1
print(random.choice(["rojo", "verde", "azul"]))  # elige un elemento al azar
print(random.shuffle)          # mezcla una lista en su lugar

lista = [1, 2, 3, 4, 5]
random.shuffle(lista)
print(lista)                   # lista mezclada
\end{lstlisting}

\begin{cajaRecuerda}
\texttt{random.randint(a, b)} incluye ambos extremos: \texttt{randint(1, 6)} puede devolver 1 o 6.
\end{cajaRecuerda}
\end{cajaEjemplo}

\subsection{Ejercicios propuestos}

\begin{cajaEjercicio}
\textbf{Ejercicio 46} \facil

Escribe un programa que use \texttt{math} para calcular y mostrar:
\begin{itemize}
    \item La raíz cuadrada de 289.
    \item El valor de $\pi$ redondeado a 4 decimales.
    \item El factorial de 10.
\end{itemize}
\end{cajaEjercicio}

\begin{cajaEjercicio}
\textbf{Ejercicio 47} \facil

Escribe un programa que simule el lanzamiento de un dado: genera un número aleatorio entre 1 y 6 y lo muestra. Pide al usuario presionar \texttt{Enter} para lanzar, y que el programa lance el dado 3 veces mostrando cada resultado.
\end{cajaEjercicio}

\begin{cajaEjercicio}
\textbf{Ejercicio 48} \intermedio

Escribe el juego \textbf{``Adivina el número''}:
\begin{enumerate}
    \item El programa elige un número aleatorio entre 1 y 50.
    \item El usuario intenta adivinarlo.
    \item Si el número es menor o mayor, el programa da una pista.
    \item Cuando acierta, muestra en cuántos intentos lo logró.
\end{enumerate}
Usa un bucle \texttt{while} que se repita hasta acertar.
\end{cajaEjercicio}

\begin{cajaEjercicio}
\textbf{Ejercicio 49} \intermedio

Escribe un programa que use \texttt{math} para calcular la \textbf{hipotenusa} de un triángulo rectángulo. Pide los dos catetos y usa la fórmula $c = \sqrt{a^2 + b^2}$. Ayuda: usa \texttt{math.sqrt()} y \texttt{math.pow()} (o \texttt{** 2}).
\end{cajaEjercicio}

\begin{cajaEjercicio}
\textbf{Ejercicio 50} \intermedio

Escribe un programa que simule \texttt{random.choice} para elegir al azar la respuesta a una pregunta. Define una lista de respuestas como \texttt{["Sí", "No", "Quizás", "Sin duda"]} y muestra una cada vez que el usuario haga una pregunta.
\end{cajaEjercicio}

\begin{cajaReto}
\textbf{Reto 11} \dificil

Escribe el juego \textbf{``Piedra, papel o tijera''} contra la computadora:
\begin{enumerate}
    \item El usuario elige piedra, papel o tijera.
    \item La computadora elige una al azar con \texttt{random.choice}.
    \item El programa determina quién gana (piedra vence a tijera, tijera vence a papel, papel vence a piedra).
    \item Se juegan 5 rondas y al final se muestra el marcador.
\end{enumerate}
Ayuda: usa \texttt{if} y \texttt{elif} para todas las combinaciones posibles.
\end{cajaReto}

% ============================================================
% MÓDULO 12
% ============================================================
\section{Proyecto Final}

El proyecto final integra todo lo aprendido. Cada estudiante debe entregar \textbf{un} programa completo, bien documentado, que funcione sin errores y que combine al menos tres temas del curso.

\subsection{Propuesta de proyectos}

\begin{cajaEjercicio}
\textbf{Proyecto A: Agenda de contactos} \intermedio

Crea un programa que gestione una agenda de contactos con un diccionario. Debe permitir:
\begin{itemize}
    \item Agregar contactos (nombre y número de teléfono).
    \item Buscar contactos.
    \item Eliminar contactos.
    \item Mostrar todos los contactos ordenados.
    \item Salir.
\end{itemize}
Temas que usa: diccionarios, listas, funciones, bucles y condicionales.
\end{cajaEjercicio}

\begin{cajaEjercicio}
\textbf{Proyecto B: Conversor de unidades} \intermedio

Crea un conversor que permita elegir entre varias conversiones (por ejemplo: kilómetros a millas, grados Celsius a Fahrenheit, kilos a libras, litros a galones). El menú se repite hasta que el usuario elija salir.
Temas que usa: funciones, condicionales, bucles, entrada de datos.
\end{cajaEjercicio}

\begin{cajaEjercicio}
\textbf{Proyecto C: Quiz de cultura general} \intermedio

Crea un cuestionario de 10 preguntas con 4 opciones cada una. El programa:
\begin{itemize}
    \item Muestra cada pregunta y las opciones.
    \item Lee la respuesta del usuario.
    \item Cuenta los aciertos.
    \item Al final muestra el puntaje y una calificación.
\end{itemize}
Temas que usa: listas, diccionarios, bucles, condicionales, entrada de datos.
\end{cajaEjercicio}

\begin{cajaEjercicio}
\textbf{Proyecto D: Generador de contraseñas} \dificil

Crea un programa que genere contraseñas seguras. Debe:
\begin{itemize}
    \item Preguntar la longitud deseada (por ejemplo, 12 caracteres).
    \item Incluir letras mayúsculas, minúsculas, dígitos y símbolos.
    \item Mezclar los caracteres al azar con \texttt{random.shuffle}.
    \item Mostrar la contraseña generada.
\end{itemize}
Temas que usa: módulo \texttt{random}, listas, cadenas, bucles, funciones.
\end{cajaEjercicio}

\begin{cajaEjercicio}
\textbf{Proyecto E: Registro de notas} \dificil

Crea un programa que lleve el registro de las notas de varios estudiantes. Debe:
\begin{itemize}
    \item Permitir registrar estudiantes con sus notas.
    \item Permitir ver las notas de un estudiante.
    \item Calcular el promedio del curso.
    \item Mostrar un reporte con el promedio de cada estudiante.
\end{itemize}
Temas que usa: diccionarios, listas, funciones, bucles, condicionales.
\end{cajaEjercicio}

\subsection{Criterios de evaluación}

El proyecto se evalúa con los siguientes criterios:

\begin{table}[h]
\centering
\begin{tabularx}{\textwidth}{l X}
\toprule
\textbf{Criterio} & \textbf{Descripción} \\
\midrule
Funcionalidad & El programa hace lo que se pide y no se cae con errores \\
Claridad & El código es legible y usa nombres descriptivos \\
Comentarios & El código tiene comentarios que explican las partes importantes \\
Uso de conceptos & Combina al menos tres temas del curso correctamente \\
Manejo de errores & El programa no falla con entradas inesperadas \\
Creatividad & Aporta ideas o funciones adicionales a lo pedido \\
\bottomrule
\end{tabularx}
\end{table}

\section{Referencias}

\begin{itemize}
    \item Documentación oficial de Python: \texttt{https://docs.python.org/es/3/}
    \item Página de descargas de Python: \texttt{https://www.python.org/downloads/}
    \item Python Tutor (visualizador de código): \texttt{https://pythontutor.com/}
    \item Replit (editor en línea): \texttt{https://replit.com/}
    \item Google Colab: \texttt{https://colab.research.google.com/}
\end{itemize}

\end{document}
